\chapter{Listen Before You Build}

Elliot Hill says he sold sneakers and T-shirts for a living. The modest
description covers a career at Nike that moved through sales, retail, and
product. When he explains how a company should grow, he puts the work in a
particular order: decide whom to serve, listen for the problem, turn the insight
into a product, tell a story that helps people recognize its value, and ask for
the order.

The order matters. A founder's imagination still has work to do, but it does
not go first and demand applause later. It enters a life already in progress.
People have routines, budgets, loyalties, fears, workarounds, colleagues, and
reasons for leaving a problem alone. Listening is how the builder discovers
that surrounding system before spending heavily to change it.

This is harder than asking people what they want. People are often generous
with encouragement and cautious with commitment. They forget ordinary
friction, predict their own behavior poorly, and describe a preferred future
that conflicts with what they did yesterday. A customer can identify a burden
without knowing the best design. A complaint can be sincere and still belong
to a tiny or unprofitable segment. A trend can be visible precisely because the
easy opportunity has passed.

Listening therefore needs a method. Its purpose is not to collect approval. It
is to let reality alter the plan while alteration is still cheap.

\section{Enter without a solution to defend}

John Abraham describes the beginning of a technology company before there is a
product to demonstrate. He starts with the decision-makers in the intended
market. What are they struggling with? What does the problem interrupt? What
would change for the company if it disappeared? Only after those conversations
does he ask whether someone will commit to a pilot and justify building a
minimum viable product.

The sequence is stronger when Chapter 5's four roles are kept separate. A
decision-maker may know the budget but not the daily workaround. A user may
know every irritation but not the contractual obstacle. The affected person
may never touch the product. The payer may value an outcome that creates more
work for everyone else. Discovery must cross the whole decision, not merely
accumulate interviews with the easiest role to reach.

Begin with events, not preferences. A useful conversation moves backward from
the last occurrence:

\begin{enumerate}[itemsep=0.2em]
\item What happened, and what triggered it?
\item What did you do next, in order?
\item Which people, tools, approvals, and payments became involved?
\item Where did time, money, attention, trust, or quality disappear?
\item What have you already tried, and why was it abandoned or retained?
\item What would have to be true for you to change the current arrangement?
\end{enumerate}

These questions leave room for a person's language and still seek observable
facts. The interviewer should resist finishing sentences, teaching the answer,
or introducing a feature merely to hear praise for it. Silence is useful. So is
the apparently irrelevant detail that appears twice in different conversations.

Abraham's final question---what would solving this mean for the company?---does
something else. It moves from irritation to consequence. The answer may reveal
that a noisy complaint has little economic weight, or that a minor delay blocks
an important sale, inspection, discharge, shipment, or renewal. This is the
outcome the later product must earn.

\section{The workaround is the first prototype}

Most problems already have a solution. It may be bad, expensive, manual,
socially awkward, or so familiar that nobody calls it a solution. A team copies
data between systems. A family keeps cash idle because the alternatives feel
opaque. A manager hires another coordinator. A patient waits. A buyer accepts
an annual inconvenience to avoid switching. These behaviors contain more
design information than a feature wish list.

Map the workaround as a sequence. Note where the person pauses, creates a
second record, asks permission, checks an answer, changes channel, or recruits
someone with informal knowledge. Record the part that must remain unchanged as
carefully as the part that fails. A replacement that saves ten minutes and
destroys a trusted approval path has made the wrong trade.

The current substitute also sets the real competition. A new tool may not be
competing with another tool. It may be competing with a spreadsheet, a trusted
employee, an insurer's rule, a habit, the possibility of doing nothing, or the
fear of being blamed for a failed change. If endurance remains cheaper and
safer than adoption, the founder has not yet found a compelling offer.

Observation needs consent and restraint. Watching work does not grant a right
to collect sensitive data, expose employees, record patients, or turn a coping
practice into marketing theater. Ask permission, minimize what is retained,
and let the people carrying the process explain what an outsider cannot see.

\section{Notice what nobody thinks to report}

A salon founder in Beverly Hills gives listening an unexpectedly physical
form. She asks employees and customers what works and what does not, but she
also watches the environment they have stopped describing. Is the front desk
welcoming? Does service move without unexplained waiting? Is the bathroom
clean? Does a chair scrape so loudly across the floor that every customer in
the room hears it?

None of these details is a grand strategy. Together they determine whether a
place feels cared for. A customer may leave with a vague sense that another
business is easier or more comfortable without ever filing a complaint about
the chair. An employee who compensates for a broken process every day may no
longer recognize the extra work as a defect. Some of the most useful evidence
is therefore silent.

Follow the complete experience at its ordinary speed. Include arrival,
orientation, waiting, handoffs, payment, error recovery, departure, and the
work employees do after the customer is gone. Look for repeated apology,
private notes, unofficial shortcuts, duplicated checks, queues with no owner,
and moments when a person must ask what happens next. Ask why a workaround
exists before removing it; an awkward step may be carrying a safety check or a
human courtesy that the formal process forgot.

Frontline workers often see these failures before managers do. Listening to
them is not a suggestion box administered from a distance. It requires enough
trust that reporting a defect will not be treated as disloyalty, and enough
authority that someone can test a repair. A company that asks for candor and
punishes the first uncomfortable answer has trained people to protect the plan
from reality.

\section{Read complaints in pairs}

An Austin e-commerce operator named Josh describes a simple spreadsheet. He
groups one- through three-star reviews on one side and four- and five-star
reviews on the other. He studies what dissatisfied buyers repeatedly dislike,
but he also preserves what satisfied buyers already value.

That second column prevents a common repair error. A company reacts to its
loudest complaint, removes the feature that loyal customers came for, and
discovers that improvement is not one-dimensional. Negative reviews identify
friction; positive reviews identify the promise that must survive the repair.

Reviews are useful because they often follow an actual purchase. They remain a
biased sample. Extreme experiences are more likely to be posted. Competitors,
incentives, moderation, product changes, and selection into the review platform
can distort the record. A five-star average can hide a weak fit for the segment
the founder hopes to serve. A detailed one-star review can describe a misuse no
responsible product should accommodate.

Sort complaints by recurrence and context before severity of language. For
each item, record the buyer type, use case, product version, date, claimed
consequence, current remedy, and whether the reviewer took a costly action such
as returning, cancelling, switching, or seeking support. Then compare the
pattern with interviews, support records, lost deals, renewals, and direct
observation.

The aim is triangulation:

\begin{equation}
\begin{gathered}
\text{credible problem evidence}\\
=\text{reported experience}+\text{observed behavior}\\
+\text{costly response}.
\end{gathered}
\end{equation}

This is a discipline, not a literal scoring formula. One exceptionally severe
failure can matter without recurrence. A common complaint can be cheap to
ignore. The equation asks the builder to seek more than one kind of evidence
before redesigning the company around a comment.

\section{Ask for a commitment that can disappoint you}

Compliments preserve relationships. Commitments expose priorities. Abraham
therefore places an early pilot before the full build. The prospective customer
must contribute something scarce: access to a workflow, staff time, data under
proper safeguards, an introduction to the payer, a signed pilot, a deposit, or
payment.

These signals form an evidence ladder:

\begin{center}
\begin{tabular}{@{}>{\raggedright\arraybackslash}p{0.31\textwidth}%
                    >{\raggedright\arraybackslash}p{0.56\textwidth}@{}}
\toprule
\textbf{Signal} & \textbf{What it can establish} \\
\midrule
Polite interest & The idea is understandable enough to discuss. \\
Past behavior & The problem occurred and produced a real response. \\
Time or access & The person will incur a small cost to explore a remedy. \\
Pilot commitment & The organization will expose the idea to operating reality. \\
Payment and use & The result competes with other uses of money and attention. \\
Renewal or referral & Value persisted beyond novelty for at least this customer. \\
\bottomrule
\end{tabular}
\end{center}

No rung proves a universal market. A deposit can be refunded. A pilot can be
sponsored by an enthusiast without authority. Payment can come from curiosity,
procurement inertia, or a temporary budget. Renewal can conceal switching
costs. The ladder is useful because each step makes agreement more expensive
and therefore less likely to be pure courtesy.

The minimum viable product belongs here, after the first commitment. Minimum
does not mean careless. A prototype for entertainment can tolerate rough
edges that a medical, financial, legal, safety, privacy, or payroll system
cannot. The smallest credible proof must preserve the standard of care while
removing every element not needed to test the uncertain mechanism.

Perfectionism sometimes protects craft. It can also protect the builder from
judgment. Abraham admits the latter possibility: a product hidden from users
cannot produce the feedback needed to correct it. The remedy is not to release
harm. It is to choose a narrower test whose failure is safe, legible, and cheap.

\section{Listen for change without worshipping trends}

Customer discovery studies the present, while a business has to survive the
future. Several speakers claim to have seen change early, and together they
show how unreliable foresight can be.

Sergio entered internet advertising when the field was young enough, in his
account, that established players had not accumulated much more knowledge than
new entrants. His earlier work selling Yellow Pages and helping with an online
program gave him a direct view of value moving from one medium to another. The
lesson is not that every new technology creates a level field. It is that a
person standing inside an old process may notice customers, money, or attention
changing direction before the new category receives a settled name.

An unidentified engineer, speaking in 2022, predicted an enormous wave of
green-technology companies. His premises were a growing world economy, heavy
dependence on fossil energy, and a long transition toward lower emissions. The
headline forecast cannot be treated as a count of future winners. The durable
question lies underneath it: which physical, regulatory, and operating
constraints must be solved even if fashion changes?

A Dubai venture and technology operator offers the opposite method: imagine
the world five years ahead and build what it will require. He names artificial
intelligence, automation, batteries, synthetic data, and other frontier
possibilities. Yet when asked about a profitable investment in the present, he
answers with mortgage technology and calls boring businesses attractive. The
contrast is valuable. A distant future can widen the search; current invoices
still reveal where demand exists now.

Test a transition thesis with four links:

\begin{equation}
\begin{aligned}
\text{durable constraint}&\longrightarrow \text{forced change}\\
&\longrightarrow \text{specific buyer burden}\\
&\longrightarrow \text{reachable first use}.
\end{aligned}
\end{equation}

If one link is missing, the forecast is still a theme. It may inspire research,
but it is not yet a market. A founder should also write the evidence that would
disprove each link: a policy reversal, a cheaper substitute, a delayed budget,
an engineering limit, customer resistance, or a distribution barrier. The
future deserves imagination and an escape route from being wrong.

\section{Listening is not obedience}

Customers possess local knowledge. They know where work hurts, which promise
was broken, what language makes sense, and what risk they fear. They may not
know which technical architecture, workflow, contract, or design can resolve
the tension across all users. Their first requested feature may preserve the
very workaround that causes the problem.

The builder's judgment is to synthesize, not vote. Preserve the need beneath a
request, compare it across segments, and test more than one response. When
evidence conflicts, choose explicitly whom the product will serve and whom it
will not. A clear boundary is more honest than a product that promises every
customer a different future.

Listening also obliges the seller to disclose information the buyer would want
to know. Later in the same Austin round of conversations, an insurance seller
presents his ability to compare products and budgets as service. That claim
becomes real only if comparison includes costs, commissions, exclusions,
guarantees, non-guaranteed projections, lapse or surrender risk, and conflicts
that could change the buyer's decision. Hearing a customer's desire while
withholding the offer's limits is targeting, not listening.

\section{Run the learning sequence}

The one-page problem brief from Chapter 5 can now become a low-cost experiment:

\begin{enumerate}[itemsep=0.1em,topsep=0.2em]
\item \textbf{Observe.} Watch the problem occur in its real setting, with
  permission, and write the sequence without proposing a feature.
\item \textbf{Interview.} Speak separately with affected people, users,
  decision-makers, payers, blockers, and former users. Ask about the last event.
\item \textbf{Map the substitute.} Record its cost, trusted elements, failure
  points, switching burden, and reason for survival.
\item \textbf{Triangulate.} Compare reported experience with reviews, support
  records, lost deals, observed behavior, and costly actions.
\item \textbf{Ask for commitment.} Seek the smallest honest contribution of
  access, time, introduction, pilot participation, deposit, or payment.
\item \textbf{Build the credible proof.} Test one uncertain mechanism at the
  lowest safe scope, then observe use rather than defending the design.
\item \textbf{Name disconfirming evidence.} Decide in advance which result will
  revise the problem, change the segment, stop the test, or end the project.
\end{enumerate}

\section{Keep a learning record}

Discovery produces conversations that are easy to remember selectively. Once
the product succeeds, every early clue can appear obvious; once it fails, every
warning can appear decisive. A short record made before the outcome protects
the company from both stories.

For each test, keep one page with seven fields:

\begin{itemize}[leftmargin=1.5em,itemsep=0.15em]
\item \textbf{Hypothesis.} The person, behavior, burden, and change believed to
  matter.
\item \textbf{Evidence.} Dated observations and conversations, with the speaker's role
  and the setting in which the behavior occurred.
\item \textbf{Contrary evidence.} The strongest fact that does not fit the preferred
  explanation, including a customer who succeeds without the product.
\item \textbf{Commitment.} What the prospective customer actually contributed and which
  authority that contribution represents.
\item \textbf{Result.} What happened in use, including adoption, failure, workaround,
  support burden, and unintended effects.
\item \textbf{Decision.} Continue, revise, narrow, change the segment, or stop, with the
  reason stated before beginning another test.
\item \textbf{Review date.} The next moment at which new evidence can change the
  decision rather than merely extend it.
\end{itemize}

The record need not preserve personal or confidential material. It should use
the minimum detail required to distinguish fact from interpretation and one
customer from a market. Record a result that changes nothing as carefully as a
breakthrough. Most learning is a narrowing of possibilities, not a dramatic
revelation.

This sequence does not remove uncertainty. It purchases better uncertainty in
small installments. At the end, the founder should be able to name the person,
old behavior, burden, decision path, tested change, and evidence of action.

A credible useful result is not yet a business. The next question is which
outcome is being bought, how much the offer can honestly claim, and what it is
worth.
